% ════════════════════════════════════════════
% 1. 项目概述
% ════════════════════════════════════════════
\section{项目概述}

\subsection{研究对象与数据}

数据取自某化妆品品牌主页的一次完整抓取，共 500 条帖子、8 个字段，无帖子编号与发布时间戳，同一主页只有一次观测，因此是帖级截面数据而非面板或时间序列。7 个字段描述发帖决策与账号状态，1 个字段记录结果：$X_1$ 为主页粉丝规模（人），$X_2$ 帖子类型（Photo / Status / Link / Video），$X_3$ 内容主题（1 / 2 / 3），$X_4$ 发布月份（1--12），$X_5$ 发布星期（1--7），$X_6$ 发布小时（1--23），$X_7$ 是否付费推广（0 / 1），$Y_1$ 为曝光总次数。表\ref{tab:head}给出前 10 条记录，量纲差异明显：粉丝规模以十万计、曝光以千至百万计，直接在原始尺度建模会放大极端值的影响。

{\small
\begin{longtable}{>{\centering\arraybackslash}p{0.130\textwidth}>{\centering\arraybackslash}p{0.110\textwidth}>{\centering\arraybackslash}p{0.100\textwidth}>{\centering\arraybackslash}p{0.080\textwidth}>{\centering\arraybackslash}p{0.080\textwidth}>{\centering\arraybackslash}p{0.080\textwidth}>{\centering\arraybackslash}p{0.100\textwidth}>{\centering\arraybackslash}p{0.160\textwidth}}
  \caption{原始数据前 10 行（列名加注中文简称）}
  \label{tab:head} \\
  \toprule
  X1 粉丝 & X2 类型 & X3 主题 & X4 月 & X5 周 & X6 时 & X7 付费 & Y1 曝光 \\
  \midrule
  \endfirsthead
  \caption{原始数据前 10 行（列名加注中文简称）（续）} \\
  \toprule
  X1 粉丝 & X2 类型 & X3 主题 & X4 月 & X5 周 & X6 时 & X7 付费 & Y1 曝光 \\
  \midrule
  \endhead
  \bottomrule
  \endfoot
  139,441 & Photo & 2 & 12 & 4 & 3 & 0 & 5,091 \\
  139,441 & Status & 2 & 12 & 3 & 10 & 0 & 19,057 \\
  139,441 & Photo & 3 & 12 & 3 & 3 & 0 & 4,373 \\
  139,441 & Photo & 2 & 12 & 2 & 10 & 1 & 87,991 \\
  139,441 & Photo & 2 & 12 & 2 & 3 & 0 & 13,594 \\
  139,441 & Status & 2 & 12 & 1 & 9 & 0 & 20,849 \\
  139,441 & Photo & 3 & 12 & 1 & 3 & 1 & 19,479 \\
  139,441 & Photo & 3 & 12 & 7 & 9 & 1 & 24,137 \\
  139,441 & Status & 2 & 12 & 7 & 3 & 0 & 22,538 \\
  139,441 & Photo & 3 & 12 & 6 & 10 & 0 & 8,668 \\
\end{longtable}
}


\subsection{分析任务}

题面要求分为四项，本文按“能不能用、怎么整理、怎么估计、怎么解释”的顺序组织：判断数据的单位、粒度、结构与各变量的因果角色，据此界定模型能回答什么；处理缺失、异常与编码变换，把原始表整理为可回归的分析数据；以曝光为因变量建立回归模型，完成共线、残差、拟合优度与预测能力诊断；在最终模型上测算付费、内容类型与内容主题的效应，并以多种设定检验结论是否依赖单一模型。

\subsection{主要结论}

四项任务的核心结果如下，具体数值见后续各章：粉丝规模与发布月份是同一条时间维度的两种写法，同时进入模型会产生不可识别的系数，模型必须以月份固定效应吸收规模趋势；在控制发帖决策变量后，付费帖子的曝光比自然帖子高 40.1\%，视频帖是图片帖的 4.90 倍，主题 2 与主题 3 分别比主题 1 低 38.0\% 与 27.0\%；付费与视频效应在六种模型设定下方向一致、量级接近，唯一需要标注样本限制的是视频帖仅 7 条；内容主题内部的付费效应并不一致，主题 1 的付费回报明显高于主题 3。
